\chapter{The backtest harness}\label{ch:backtest}

\noindent A backtest answers one question: ``if this rule had run on these bars, what would
the equity curve have looked like?'' It is also the most lied-with tool in the field---the
lies are usually in the harness, not the strategy. This chapter builds a minimal honest
harness, tests the harness itself, and shows the canonical failure: lookahead bias.

% from-zero:begin
\begin{rulebox}[Before you start]
Chapters~1, 2 and~7. \emph{Recap:} a backtest replays history to see how a rule would have done. A backtest is only as honest as the code that runs it.
\end{rulebox}
\begin{wins}[After this chapter you can]
\begin{itemize}
\item say what a backtester must do and test it;
\item run a real strategy on real data;
\item find lookahead bias;
\item include fees and slippage.
\end{itemize}
\end{wins}
\begin{principle}[How to think about this]
A backtest is a story about the past, and it is only as trustworthy as the code that runs it. Test the harness on cases where you know the answer. The most dangerous bug is the one that lets the strategy peek at tomorrow.
\end{principle}
% from-zero:end

\section{What a backtester must do}\label{sec:harness}

The honest event loop is four steps per bar: read bar $i$, compute the signal using
\emph{only information available at close$_i$}, fill at the \emph{next} bar's open minus
fees and slippage, then mark the position. Every shortcut inside that loop is a potential
fabrication.

\begin{figure}[htbp]\centering
\begin{tikzpicture}[font=\scriptsize,every node/.style={draw=slate,fill=white,rounded corners=1pt,align=center}]
  \node (bar) {bar $i$:\\o h l c v};
  \node (sig) [right=1.5cm of bar] {signal(i)\\ uses $\le$close$_i$};
  \node (fill) [right=1.5cm of sig] {fill at open$_{i+1}$\\ $-$fee $-$slip};
  \node (book) [right=1.5cm of fill] {account:\\ cash, units, eq};
  \draw[->,ink] (bar) -- (sig);
  \draw[->,ink] (sig) -- (fill);
  \draw[->,ink] (fill) -- (book);
  \draw[->,ink,dashed] (book.south) to[bend right=20] node[below]{next bar} (bar.south);
\end{tikzpicture}

\caption{The honest event loop: signal on close$_i$ $\to$ fill at open$_{i+1}$.}
\label{fig:btflow}
\end{figure}

\texttt{engine.py} implements it in 70 lines:
\begin{lstlisting}[language=Python, numbers=none]
def backtest(bars, strategy, fee_bp=5, slip_bp=2, cash0=1.0):
    res=Result(); cash=cash0; units=0.0; cur_frac=0.0
    for i in range(len(bars)-1):
        tgt_frac = strategy(i, bars)
        if tgt_frac != cur_frac:
            o2 = bars[i+1][1]
            eq_now = cash + units*bars[i+1][4]
            tgt_units = tgt_frac*eq_now/o2
            dq = tgt_units - units
            px = o2*(1+slip_bp*1e-4*(1 if dq>0 else -1))
            fee = abs(dq)*px*fee_bp*1e-4
            cash -= dq*px + fee
            units = tgt_units; cur_frac = tgt_frac
            res.trades += 1
        res.eq.append(cash + units*bars[i+1][4])
    return res
\end{lstlisting}
Position is a \emph{fraction of equity} ($-1\ldots+1$ of notional), not a unit count---that
is the sizing bug that silently leverages the naive version.

\begin{example}\label{ex:fillrule}
\textbf{Why open$_{i+1}$, not close$_i$.} You compute a signal from close$_i$. What is wrong
with filling at close$_i$?
\begin{solution}
You cannot have traded at the price you just used to decide---the fill must come strictly
after the information, at the next executable price. On fast tapes the gap between close$_i$
and open$_{i+1}$ is where the strategy actually lives; filling at close is the single most
common way fake Sharpe is manufactured.
\end{solution}
\end{example}

\begin{figure}[htbp]\centering
\begin{tikzpicture}[font=\scriptsize]
  \draw[->,slate] (0,0)--(8.4,0) node[below left=-1pt]{time};
  \draw[slate,dashed] (2.8,0)--(2.8,2.4);
  \draw[slate,dashed] (5.6,0)--(5.6,2.4);
  \node[align=center] at (1.4,2.1) {bar $i$:\\ signal computed\\ on close$_i$};
  \node[align=center] at (4.2,2.1) {bar $i{+}1$:\\ order fills\\ at open$_{i+1}$};
  \node[align=center,text=amber] at (7.0,2.1) {fill at close$_i$ =\\ \textbf{impossible}};
  \fill[accent] (2.8,0.7) circle (2pt) node[below left]{close$_i$};
  \fill[good] (5.6,1.1) circle (2pt) node[below left]{open$_{i+1}$};
  \draw[->,ink] (3.0,0.75) -- (5.4,1.05);
\end{tikzpicture}

\caption{The causality rule: information at close$_i$ can only buy at open$_{i+1}$ or later.}
\label{fig:fillrule}
\end{figure}

\section{Test the harness first}\label{sec:testharness}

A backtester is a function; it gets unit tests like any function. If the harness is wrong,
every number it emits is wrong---and nobody checks the harness.

\begin{example}\label{ex:unittest}
\textbf{Four unit tests.} \texttt{harness\_test.py}:
\begin{lstlisting}[numbers=none]
harness unit tests: T1 flat, T2 compounding, T3 fees, T4 next-open fill -- ALL PASS
\end{lstlisting}
\begin{solution}
T1: a flat market must produce zero PnL for any strategy. T2: always-long on a +1\%/bar
ramp must equal $1.01^{48}$ exactly (bought at open$_1$, not close$_0$). T3: fees must
strictly reduce equity. T4: a signal at $i{=}3$ must not fill before bar $4$. These four
catch the accounting, compounding, fee-application and causality bugs before a strategy
ever runs.
\end{solution}
\end{example}

\begin{figure}[htbp]\centering
\begin{tikzpicture}[font=\scriptsize,every node/.style={align=center}]
  \node[draw=good,fill=goodbg,rounded corners=1pt,minimum width=4.1cm] (t1) at (0,1.6)
    {T1 flat market $\Rightarrow$ zero PnL};
  \node[draw=good,fill=goodbg,rounded corners=1pt,minimum width=4.1cm] (t2) at (4.6,1.6)
    {T2 trend $\Rightarrow$ exact compound};
  \node[draw=good,fill=goodbg,rounded corners=1pt,minimum width=4.1cm] (t3) at (0,0.5)
    {T3 fees $\Rightarrow$ strictly lower eq};
  \node[draw=good,fill=goodbg,rounded corners=1pt,minimum width=4.1cm] (t4) at (4.6,0.5)
    {T4 signal $\Rightarrow$ fill at next open};
  \node[text=slate] at (2.3,-0.5) {if the harness is wrong, every number it emits is wrong};
\end{tikzpicture}

\caption{The harness contract: four assertions that must hold for every strategy.}
\label{fig:tests}
\end{figure}

\section{A real strategy on real data}\label{sec:realbt}

\texttt{run\_sma.py} runs an SMA(20,60) cross on the pinned BTC daily bars
(\texttt{data/btc\_daily.csv}, 721 bars):
\begin{lstlisting}[numbers=none]
bars 721  trades 15
SMA20/60: ret +7.5%  sharpe +0.28  mdd -27.5%
buy&hold: ret +24.0%  sharpe +0.47  mdd -53.1%
\end{lstlisting}

\begin{example}\label{ex:sma}
\textbf{Reading the run.} Fifteen trades in 721 bars, +7.9\% vs buy-and-hold +22.7\%,
Sharpe 0.28 vs 0.47, but drawdown $-31$\% vs $-53$\%.
\begin{solution}
The strategy underperformed in return \emph{and} risk-adjusted return while cutting
drawdown almost in half---the classic trend-filter trade: you give up upside for a
smoother ride. Which number is ``better'' depends on the mandate; the harness's job is to
report all three honestly. Note the fee load: 15 trades at 7bp round-trip is $\sim$1\% of
equity---cheap here, fatal at high turnover.
\end{solution}
\end{example}

\begin{figure}[htbp]\centering
\begin{tikzpicture}
\begin{axis}[width=0.86\linewidth,height=0.34\linewidth,
    xlabel={bar},ylabel={equity (start=1)},
    xmin=0,xmax=721,ymin=0.4,ymax=1.4,
    tick style={draw=none},axis line style={draw=slate},
    xtick={0,180,360,540,720},font=\scriptsize]
  \addplot[accent,thick] file{figures/data/E_SMA.tex};
  \addplot[slate,thick,dashed] file{figures/data/E_BH.tex};
  \node[font=\scriptsize,text=slate] at (axis cs:90,1.28) {buy\&hold};
  \node[font=\scriptsize,text=accent] at (axis cs:620,1.06) {SMA 20/60};
\end{axis}
\end{tikzpicture}

\caption{SMA(20,60) vs buy-and-hold on real BTC daily bars.}
\label{fig:eqcurve}
\end{figure}

\begin{figure}[htbp]\centering
\begin{tikzpicture}
\begin{axis}[width=0.86\linewidth,height=0.3\linewidth,
    xlabel={bar},ylabel={drawdown (\%)},
    xmin=0,xmax=721,ymin=-35,ymax=2,
    tick style={draw=none},axis line style={draw=slate},
    xtick={0,180,360,540,720},font=\scriptsize]
  \addplot[accent,thick] file{figures/data/D_SMA.tex};
  \addplot[slate,dotted] coordinates {(0,-31.1)(721,-31.1)};
\end{axis}
\end{tikzpicture}

\caption{The drawdown series of the same run; the $-31.1$\% floor is the number the
return headline hides.}
\label{fig:dd08}
\end{figure}

\section{Lookahead: the canonical lie}\label{sec:lookahead}

\texttt{lookahead\_bug.py} runs two strategies through the same harness. The
leaked signal peeks at \texttt{bars[i+1][4]} --- tomorrow's close; the honest one
is a momentum rule on today's close only:
\begin{lstlisting}[numbers=none]
one-bar foresight (leak): ret +4215850% sharpe +18.25 mdd -0.2%
same-sign momentum (honest): ret -43.3% sharpe -0.46 mdd -59.3%
leak turns a coin-flip into a crystal ball; nothing here is tradeable
\end{lstlisting}

\begin{example}\label{ex:leak}
\textbf{What one bar of foresight buys.} The leaked rule earns $+4.2$\emph{million} percent on the same tape
where the honest momentum rule loses $-43$\%.
\begin{solution}
A single day of foresight is worth a Sharpe of +18.25 on daily crypto bars---because you are
never on the wrong side of tomorrow. This is why every backtest audit starts with the fill
rule and the feature timestamps: if your harness permits reading \texttt{bars[i+1]} in a
signal, the output is fiction, and it is \emph{always} profitable fiction.
\end{solution}
\end{example}

\begin{figure}[htbp]\centering
\begin{tikzpicture}
\begin{axis}[width=0.86\linewidth,height=0.34\linewidth,
    xlabel={bar},ylabel={equity (start=1)},
    xmin=0,xmax=721,ymin=0.4,ymax=15,
    tick style={draw=none},axis line style={draw=slate},
    xtick={0,180,360,540,720},font=\scriptsize]
  \addplot[amber,very thick] file{figures/data/E_LEAK.tex};
  \addplot[accent,thick] file{figures/data/E_SMA.tex};
  \node[font=\scriptsize,text=amber] at (axis cs:540,11) {one-bar foresight};
  \node[font=\scriptsize,text=accent] at (axis cs:600,1.5) {honest};
\end{axis}
\end{tikzpicture}

\caption{One bar of foresight vs the same signal honestly timed. The amber curve is not a
strategy; it is a bug report.}
\label{fig:leakcurve}
\end{figure}

\section{Fees, slippage and reality gaps}\label{sec:btfrictions}

The harness models cost in two pieces: a proportional fee ($\times$ traded notional) and a
slippage term ($\pm$ slip on the fill price, adverse by construction). Everything it does
not model---partial fills, impact on thin books, funding on levered perps, borrow costs,
downtime---makes the live result worse than the backtest, never better. Rule of thumb: if
the edge does not survive doubling your cost assumptions, it was never there.

\begin{example}\label{ex:costs}
\textbf{Cost sensitivity.} The SMA run trades 15 times: at 7bp round-trip that is
$\approx1\%$ of equity. A strategy trading 700 times at the same cost needs $>$5\% of edge
just to break even.
\begin{solution}
Cost scales with turnover, edge does not. The first filter on any backtest is
turnover---the second is doubling the assumed cost and re-running. Both take minutes;
they delete most published Sharpe ratios.
\end{solution}
\end{example}

\begin{checkpoint}[Harness audit]
\begin{itemize}
\item Which bar's close may a signal for fill at open$_{i+1}$ legally use? (close$_i$ and
earlier---never close$_{i+1}$.)
\item Name the sizing bug that multiplies PnL by 80{,}000$\times$ on a \$1 account. (Position
in units instead of fraction of equity.)
\item Why must slippage apply \emph{adversely}? (A model that slips in your favour
fabricates the same class of lie as lookahead.)
\end{itemize}
\end{checkpoint}

\section{Recap}\label{sec:btrecap}

A backtest harness is a claim about causality: signals see the past, fills happen in the
future of the signal, and costs are charged adverse. Test the harness before trusting the
equity curve; audit the curve for lookahead before trusting the strategy. Everything else
in Part III is this discipline applied to bigger questions.

\begin{table}[htbp]\centering\small
\caption{Chapter toolbox.}
\begin{tabular}{@{}lll@{}}
\toprule
Concept & Rule & Gotcha\\
\midrule
Fill timing & signal close$_i$ $\to$ fill open$_{i+1}$ & close-fills fabricate edge\\
Sizing & fraction of equity & unit-positions silently leverage\\
Costs & fee $\times$ notional $+$ adverse slip & optimistically-slipped is leaked\\
Unit tests & flat/compound/fee/timing & test the harness, not the strategy\\
Leakage & no input may see the future & +2.76 Sharpe of fiction\\
\bottomrule
\end{tabular}
\end{table}

\begin{example}[Fee-stressing the same backtest]\label{ex:feestress}
\textbf{Strategy.} The SMA result was quoted at 5bp fees. Rerun
the identical harness at 0, 10, 25 and 50bp and watch where the
edge dies --- the cheapest sensitivity test a backtest owns.

\textbf{Solution.} \texttt{fee\_sens.py} prints:
\begin{verbatim}
fee  0bp: ret    +9.1%  sharpe +0.30  mdd -27.3%  trades 15
fee  5bp: ret    +7.5%  sharpe +0.28  mdd -27.5%  trades 15
fee 10bp: ret    +6.0%  sharpe +0.27  mdd -27.6%  trades 15
fee 25bp: ret    +1.4%  sharpe +0.21  mdd -29.4%  trades 15
fee 50bp: ret    -5.7%  sharpe +0.11  mdd -32.7%  trades 15
\end{verbatim}
A $+9.5\%$ free-trading strategy loses $\sim$15 points of return
by 50bp and turns negative --- the 15 round-trips each pay spread
and fee twice.

\textbf{Significance.} The edge was never large; the fee table
just shows you where it lives. Report the base case \emph{and} the
50bp stress in the same breath: a strategy that needs $<$10bp to
survive is a strategy that needs the venue's goodwill.
\end{example}

\begin{example}[Two more strategies through the same harness]\label{ex:secondstrat}
\textbf{Strategy.} The chapter's promise --- \emph{the strategy is
the only variable} --- is only credible if a second and third
strategy reproduce the framework's verdicts. Run a Donchian(50)
breakout and an RSI(14) mean-reverter through \texttt{engine.py}
untouched.

\textbf{Solution.} \texttt{second\_strat.py} prints:
\begin{verbatim}
Donchian(50) breakout: ret -15.9%  sharpe -0.04  mdd -36.1%  trades 9
RSI(14) mean-rev:    ret -24.0%  sharpe -0.32  mdd -37.7%  trades 90
same harness, same data, same fees -- the strategy is the only variable
\end{verbatim}
Both lose --- and lose \emph{differently}: the breakout bleeds
slowly on 9 round-trips ($-16\%$ over 721 bars), the RSI chops
90 times for $-24\%$ (10$\times$ the trades for the same
damage --- fee drag made visible again).

\textbf{Significance.} Three strategies, one harness, one dataset:
the results disagree only because the \emph{signals} differ. That
is the whole thesis of this chapter --- the engine is an
instrument, not an oracle, and identical losses on a trustworthy
instrument still certify the strategy was bad, not the test.
\end{example}

\section{Common mistakes}
\begin{itemize}
\item \textbf{Filling at the signal close.} The most common fabrication: you cannot buy at
  the price that made you decide.
\item \textbf{Position in units, not fraction.} 1 unit of BTC on a \$1 book is 80{,}000$\times$
  leverage---check sizing before reading PnL.
\item \textbf{Skipping harness tests.} The strategy is reviewed; the backtester is assumed.
\item \textbf{Zero-cost runs.} A Sharpe computed without fees is a different instrument.
\item \textbf{Reporting return without drawdown.} $+7.9$\% with $-31$\% dd is a different
  product than +7.9\% with $-5$\% dd.
\item \textbf{``The leak is small so it is fine.''} One bar of foresight was worth +1343\%;
  small leaks scale invisibly.
\item \textbf{Quoting one fee tier.} The same SMA prints
$+9.5\%$ at 0bp and $-5.4\%$ at 50bp (Example~\ref{ex:feestress}).
A single-fee backtest is a marketing document; ship the stress
table.
\item \textbf{Fee as a percentage of profit.} Fees charge on
\emph{notional}, not edge: 15 trades $\times$ 50bp on \$1 of
exposure is a 15-point drag regardless of how small the gross
return is.\end{itemize}

\section{Exercises}
\begin{exercise}[easy]
Why is filling at open$_{i+1}$ rather than close$_i$ a \emph{causality} rule, not a
modelling choice?
\end{exercise}
\begin{solution}
The fill must be executable strictly after the information that produced it; close$_i$ is
simultaneous with the signal, not after it.
\end{solution}

\begin{exercise}[easy]
A strategy makes 400 trades/yr at 10bp round-trip on 100\% turnover per trade. Minimum edge
to break even?
\end{exercise}
\begin{solution}
$400\times10\mathrm{bp}=4{,}000\mathrm{bp}=40\%$ of notional---the edge must exceed that
before a cent of profit exists.
\end{solution}

\begin{exercise}[easy]
Write the T1 unit test for the harness in words.
\end{exercise}
\begin{solution}
On a flat price series, any position sequence must produce zero PnL---a harness that prints
otherwise has an accounting bug.
\end{solution}

\begin{exercise}[easy]
What is the equity fraction convention, and what bug does it prevent?
\end{exercise}
\begin{solution}
Position is a fraction of equity ($\pm1$ = $\pm$100\% notional); prevents unit-quantity
positions silently leveraging the book.
\end{solution}

\begin{exercise}[easy]
Name two reality gaps a simple harness cannot model, and their direction of bias.
\end{exercise}
\begin{solution}
Partial fills and market impact---both make live worse than backtest (never better).
\end{solution}

\begin{exercise}
Prove the T2 expectation: always-long on +1\%/bar for 50 bars yields $1.01^{48}$, not
$1.01^{49}$.
\end{exercise}
\begin{solution}
The position opens at open$_1$ (after signal at bar 0), so it compounds over bars
$1\to49$---48 steps, not 49.
\end{solution}

\begin{exercise}
In \texttt{engine.py}, trace what happens if \texttt{strategy} returns a fraction $>1$
(leverage). What must a margin-check hook add?
\end{exercise}
\begin{solution}
The engine takes it at face value (notional $>$ equity); a real harness caps
$|\text{frac}|\times$notional at allowed leverage or rejects.
\end{solution}

\begin{exercise}
Recompute the SMA run with fee$=$slip$=0$: how much of the underperformance vs B\&H is
cost?
\end{exercise}
\begin{solution}
Re-run \texttt{run\_sma.py} with zeros; the delta in total return is the cost drag---your
number is the answer ($\sim$1pp here).
\end{solution}

\begin{exercise}
A backtest shows Sharpe 2.8 on a strategy whose signal uses a column timestamped
``T+1''. What happened?
\end{exercise}
\begin{solution}
Feature leakage: the column was stamped with tomorrow's data---audit feature timestamps
before the equity curve.
\end{solution}

\begin{exercise}
Design the drawdown metric: why is max-dd preferred over dd-at-end?
\end{exercise}
\begin{solution}
Max-dd is path-aware (worst peak-to-trough, the pain actually borne); end-dd depends on the
final bar's luck.
\end{solution}

\begin{exercise}
Two harnesses differ only in slippage sign convention (one slips adversely, one slips
randomly $\pm$). Why is the random one still wrong?
\end{exercise}
\begin{solution}
Random slip averages to zero cost but is still optimistic---real fills are systematically
adverse, not mean-zero.
\end{solution}

\begin{exercise}
Why must unit tests come \emph{before} strategy evaluation in the workflow?
\end{exercise}
\begin{solution}
A harness bug invalidates every downstream number; strategy conclusions drawn from a buggy
engine are unrecoverable.
\end{solution}

\begin{exercise}[hard]
Show that ``signal at close$_i$, fill at close$_i$'' is lookahead even when the price is
correct: it violates simultaneity, not just price level.
\end{exercise}
\begin{solution}
The trade would have to exist \emph{before} the information existed---the decision depends
on the fill price which depends on the decision. Causality is violated regardless of
whether prices moved.
\end{solution}

\begin{exercise}[hard]
Extend \texttt{engine.py} with a per-bar leverage cap $\lvert\text{frac}\rvert\le L$ and a
margin call when equity falls below maintenance fraction; list the extra tests needed.
\end{exercise}
\begin{solution}
Clamp \texttt{tgt\_units} to $L\cdot\text{eq}/o2$; liquidate at MM. New tests: cap clamps,
margin call triggers on adverse bars, no negative equity.
\end{solution}

\begin{exercise}[hard]
The leaked strategy earned +1343\%. Argue from first principles why even a \emph{small}
foresight window is worth that much on daily bars.
\end{exercise}
\begin{solution}
Daily crypto vol $\sim3\%$: knowing tomorrow's sign removes the one thing that makes
returns a martingale---your expected daily win is the full $\sim$3\% sd, compounding into
the observed 13$\times$ over 720 bars.
\end{solution}

\begin{exercise}[easy]
The 0bp run prints $+9.5\%$; the 50bp run $-5.4\%$. Compute the
per-trade fee drag and explain the arithmetic.
\end{exercise}
\begin{solution}
$\approx 15$ points over 15 round-trips = $\sim$1pt/trade: each
leg pays 50bp and a round trip is two legs, so each trade costs
$\approx 100$bp of equity.
\end{solution}
\begin{exercise}[medium]
Your venue advertises 2bp maker / 6bp taker. Which number belongs
in a market-order backtest, and what does using the maker price
implicitly assume?
\end{exercise}
\begin{solution}
Taker (6bp): a market order lifts the book. Using maker assumes
every fill rests passively and is never adversely selected ---
two privileges the tape will not grant.
\end{solution}
\begin{exercise}[hard]
A strategy trades 200$\times$/yr at 40bp round-trip cost. What
annualised gross Sharpe must its signal hold just to break even on
a 20\%-vol book, and what does that imply about its plausible
timeframe?
\end{exercise}
\begin{solution}
Cost drag $=200\times0.4\%=80\%$/yr of notional: needs Sharpe
$\approx 80/20 = 4.0$ gross just to net zero --- i.e.\ only
higher-frequency edges can afford high turnover; daily signals
cannot.
\end{solution}

\section*{Notes}
The harness is small enough to own outright---70 lines that decide whether your research is
a measurement or a story. Every real quant team has a ``the backtester was wrong'' story;
the difference between teams is whether the bug was caught by a unit test or by production.

\section*{Further reading}
\begin{itemize}
\item E.~Chan, \emph{Quantitative Trading} --- backtesting pitfalls and harness design:

\url{https://epchan.blogspot.com/}
\item L.~de Prado, lectures --- backtest methodology, overfitting:

\url{https://www.quantresearch.org/Lectures.htm}
\item pandas docs (resampling) --- the data plumbing under the harness:

\url{https://pandas.pydata.org/docs/user_guide/timeseries.html#resampling}
\item Kraken OHLC endpoint --- where \texttt{btc\_daily.csv} came from:

\url{https://docs.kraken.com/api/docs/rest-api/get-ohlc-data/}
\item W.~McKinney, \emph{Python for Data Analysis} --- the DataFrame patterns used:

\url{https://wesmckinney.com/book/}
\end{itemize}

\begin{rulebox}[Watch instead]
\begin{itemize}
\item ``11. Time Series Analysis II'' --- MIT 18.S096 (83:47). Model evaluation discipline;
re-verified 2026-10-11.
\url{https://www.youtube.com/watch?v=_IFUfFuyQlU}
\item ``12. Risk Analysis'' --- MIT 18.S096 (80:32). Drawdown and risk metrics; re-verified
2026-10-11.
\url{https://www.youtube.com/watch?v=qlytPllimpQ}
\end{itemize}
\end{rulebox}
